% Physical unit 026; accepted Session 03 after fresh 400-dpi authority replay.
\GordanUnitStart{026}{printed page 21}
\GordanTopFolio{21}
\begingroup\small
vel referendo ad universum genus:
\[
 \omega+4n\mathcal A(l_0)-2h\pi
\]
qua in formula $n$ deinceps omnium integrorum numerorum valores induit, et $h$ eum numerum integrum repraesentat, qui est in:
\[
 {\psi_2+4n\mathcal A(l_0)\over 2\pi}.
\]
Quod igitur
\[
 {\mathcal A(l_0)\over 2\pi}
\]
irrationale est (quippe contrarium, quod rarissimo fit, omittendum est); $\psi$ omnes deinceps valores, intra $0$ et $2\pi$ sitos, assumit; id, quod erat demonstrandum.

\begin{center}\S\ 4.\end{center}

Jam vero, priusquam quaerere instituimus, quae sit brevissima omnium linea intra duo puncta $A(l\psi)$ et $B(l_1\psi_1)$ in sphaeroide; nonnulla eam rem spectantia considerare placeat.

Illa in quaestione $l$ majus esse $l_1$ ponamus. E theoremate superiori cum jam dilucide appareat, inter lineas a puncto $A$ exeuntes omnes punctum $B$ continere; quaerendum est accuratius, qualis ceteris brevior existat. Ubi haec quidem res in promptu est, lineam, quae ceteris sit brevior, genere noto, definiri posse.

\begin{center}ARGUMENTATIO.\end{center}

Nam si alicujus lineae geodeticae cursum a puncto $A$ atque multarum aliarum linearum geodeticarum prope hanc currentium et minus in meridianum puncti $A$ inclinatarum intuemur, lineam geodeticam a linea aliqua eo prius secari, invenimus, quo minor sit inclinatio in ipsam. Qua e re sequitur, ut omnes lineae geodeticae ejusdem generis hanc post terminum secent. Quum vero quaevis linea geodetica, ut supra demonstratum est, usque ad terminum brevissima sit; eo magis, lineam aliam geodeticam secans, jam longe linea esse brevissima desierit, necesse est.

\begin{center}\S\ 5.\end{center}

Priusquam autem ad genus determinandum propius accedamus, duo theoremata proponamus haec.

I. \gq{Linea geodetica $l_0$ I, si a puncto $A$ exorsa, nondum assecuta punctum $B$, secat meridianum puncti $A$ (e. g. in $C$); non est linea omnium brevissima inter $A$ et $B$.}

\begin{center}ARGUMENTATIO.\end{center}

Ea linea geodetica II, cujus Azimut negativo Azimut lineae geodeticae I in $A$ aequale est, propter symmetriam aeque ac linea I et tangit circulum $l_0$ et continet
\endgroup
